% 基础篇：原稿模块一至三及 E2、E4、E5、F2--F7、G1、G2。
\section{从照片到可以计算的数字图像}
\label{sec:foundations-image}

一张手机照片看上去有天空、道路和行人，计算机读到的却是一份按位置排列的数值。学习图像处理的第一步，是把这两种看法联系起来：数值记录某个位置的明暗或颜色，数值之间的关系记录物体的形状。处理图像，就是在保留任务所需信息的前提下，有目的地改变或解释这些数值。

本章从方格和数表开始，不要求预先学过信号处理。后续章节中反复出现的“逐像素”“邻域”“灰度”“分辨率”，都可以回到这里找到具体含义。

\subsection{一幅图像为什么可以写成一张数表}

把一张黑白照片盖上透明方格纸，每格用一个数代表明暗，就得到数字图像（digital image）的直观模型。方格中的一个记录单位叫像素（pixel）。像素不是被拍物体上的一个天然小方块，而是数字图像中的一个位置及其数值；把照片放大后看到的方块，是显示软件画出像素的一种方式。

对只有明暗、没有颜色的灰度图像（grayscale image），每个像素用一个灰度级（gray level）表示。常见的 8 位图像把灰度编号为 0 到 255：0 通常显示为黑，255 通常显示为白，中间的数显示为不同深浅的灰。这里的“通常”很重要：编号如何变成屏幕亮度，还要经过显示设备和显示设置；灰度 200 并不直接等于某个固定物理亮度。

\begin{figure}[htbp]
\centering
\begin{tikzpicture}[x=.78cm,y=.78cm,font=\small]
% Display white for 255 and black for 0, explicitly overpaint correct gray mapping.
\foreach \x/\y/\v in {0/2/0,1/2/60,2/2/180,3/2/255,0/1/20,1/1/90,2/1/200,3/1/240,0/0/50,1/0/120,2/0/220,3/0/255}{
  \pgfmathsetmacro{\shade}{100-\v/2.55}
  \fill[black!\shade] (\x,\y) rectangle ++(1,1);
}
\draw[step=1,gray] (0,0) grid (4,3);
\draw[-{Latex}] (0,3.5)--(4,3.5) node[right]{列 $j$};
\draw[-{Latex}] (-.45,3)--(-.45,0) node[below]{行 $i$};
\foreach \x in {0,...,3} \node at (\x+.5,3.25) {\x};
\foreach \y/\v in {2/0,1/1,0/2} \node at (-.2,\y+.5) {\v};
\draw[line width=1.2pt] (1,1) rectangle (2,2);
\draw[-{Latex}] (4.4,1.5)--(5.3,1.5);
\node[anchor=west] at (5.6,1.5) {$\displaystyle f=\begin{bmatrix}0&60&180&255\\20&90&200&240\\50&120&220&255\end{bmatrix}$};
\node at (2,-.65) {3 行、4 列，共 12 个像素};
\node[anchor=west] at (5.6,-.05) {$f(1,1)=90$：第 2 行、第 2 列};
\end{tikzpicture}
\caption{小灰度图与数表的一一对应。这里从 0 开始编号，因此编号 1 指第二行或第二列。}
\label{fig:foundation-grid}
\end{figure}

这种数表在数学中叫矩阵（matrix）。记图像为 $f(i,j)$，意思是“到第 $i$ 行、第 $j$ 列去取一个灰度数”。字母 $f$ 是图像的名字，括号里的 $i,j$ 是地址，不是乘法。若图像有 $M$ 行、$N$ 列，并从 0 编号，则
\[
  i=0,1,\ldots,M-1,\qquad j=0,1,\ldots,N-1.
\]
例如 $M=3,N=4$，共有 $MN=3\times4=12$ 个像素。教材也常写连续坐标 $f(x,y)$，其中 $x$ 向右、$y$ 向上；图像数组常把行号向下增加。两者都可以使用，但计算位置或距离前必须先说明约定。本基础篇的离散数表统一使用“行 $i$、列 $j$”。

\subsection{二值 灰度和彩色三种记录方式}

二值图像（binary image）只使用两个值，通常记为 0 和 1。例如，把车牌文字记为 1、背景记为 0，就得到“哪里是文字”的地图。这里的 1 表示被选中的目标，不一定必须显示成白色。判断像素是不是目标，与目标显示成什么颜色，是两个不同问题。

彩色图像（color image）通常为每个像素记录三个通道（channel），例如红绿蓝颜色模型（red--green--blue color model, RGB）中的 $R,G,B$。像素 $(255,0,0)$ 可显示为红色，$(0,255,0)$ 可显示为绿色，$(80,80,80)$ 可显示为灰色。一个三通道 RGB 图像可以看成三张大小相同的灰度数表叠在一起。三通道不是“一个像素只能有三种颜色”，而是用三个强度组合表达大量颜色。

记录一个二选一结果需要一个比特（bit）。每增加一个比特，可用组合数就翻倍。因此，$b$ 比特可以表示
\[
  L=2^b
\]
个不同编号。式中 $b$ 是每个灰度像素的位数，$L$ 是可用灰度级数。例如 $b=1$ 时 $L=2$，$b=8$ 时 $L=256$，$b=10$ 时 $L=1024$。8 比特又叫 1 字节（byte）。一张 $1000\times1000$、每像素 8 比特的单通道图像，未压缩像素数据为
\[
  1000\times1000\times8=8\,000\,000\ \text{bit}
  =1\,000\,000\ \text{byte}.
\]
同样尺寸、每个通道 8 比特的 RGB 图像，像素数据是它的三倍。实际文件还可能包含拍摄信息并经过压缩，不能仅凭文件大小反推像素数。

\subsection{图像处理要完成什么任务}

想看清一张昏暗的门禁照片时，输入是图像，输出通常仍是图像：把暗部提亮、减弱噪点、让边缘更清楚。这属于低级图像处理（low-level image processing）常见的任务。图像改善（image improvement）强调让人更容易观察和解释，也可以为后面的机器分析准备更合适的输入。低级并不表示技术简单或效果不重要，而是指离像素数据比较近的处理层次。

当任务变成“这张照片里有几个人”或“这个号码是什么”时，输出已经是数量、位置或类别。图像分析（image analysis）和高级图像处理（high-level image processing）更关注结构、特征与语义。其中，图像分割（image segmentation）把像素分成有意义的区域，特征提取（feature extraction）把形状、颜色或纹理概括成可比较的量，识别（recognition）进一步判断对象属于哪一类。低级处理常为这些步骤服务。本讲义原课程的重点是低级图像处理，尤其是图像增强。

这些层次之间没有一条适用于所有教材的硬边界。例如，压缩表示（compressed representation）既能为存储和传输节省空间，也能服务后续分析，不能仅因为它出现在分析流程里就认定为“只属于高级处理”。理解具体任务的输入、输出和目的，比死记分类更可靠。

\begin{figure}[htbp]
\centering
\begin{tikzpicture}[font=\small,node distance=.52cm]
\node[draw,align=center,text width=2.5cm,minimum height=1.05cm] (a) {图像获取\\把场景记录成图像};
\node[draw,align=center,text width=2.5cm,minimum height=1.05cm,right=of a] (b) {图像增强与复原\\改善可见性与质量};
\node[draw,align=center,text width=2.5cm,minimum height=1.05cm,right=of b] (c) {分割与特征提取\\找到目标并描述};
\node[draw,align=center,text width=2.5cm,minimum height=1.05cm,right=of c] (d) {识别与决策\\回答实际问题};
\draw[-{Latex}] (a)--(b);\draw[-{Latex}] (b)--(c);\draw[-{Latex}] (c)--(d);
\node[align=center,below=.55cm of b,text width=10cm,xshift=1.6cm] {彩色处理、几何校正、压缩和存储可按任务插入流程；\\实际系统不必把所有步骤依次执行一遍。};
\end{tikzpicture}
\caption{从图像到任务答案的一条典型路线。流程由用途决定。}
\end{figure}

图像处理属于信号处理（signal processing）的一部分。信号可以理解为“随位置或时间变化、携带信息的量”：录音随时间变，照片随二维位置变。它也可以作为工程工具，例如用图像判断工件是否缺角。传统课程从电路分析、信号与系统、随机信号处理逐步进入图像处理；本书先从可见的像素操作建立直觉，再补上需要的数学语言。

计算机图形学（computer graphics）主要研究如何根据模型、几何和材质生成图像，例如根据三维模型合成汽车广告；图像处理则常从已有图像出发。计算机视觉（computer vision）关心机器如何通过图像理解场景，图像处理常是它的组成步骤。模式识别（pattern recognition）研究如何依据特征判别类别，既可处理图像，也可处理语音等数据。人工智能（artificial intelligence）与图像处理、图像识别相交；“图像处理只做像素而人工智能只做识别”不是可靠的分界。

\subsection{相机以外也可以产生图像}

图像获取（image acquisition）把来自场景的测量结果组织成二维或更高维数据。普通摄像机记录可见光；鱼眼相机（fisheye camera）用很宽的视角观察环境，画面也因此有明显的几何变形；全向相机（omnidirectional camera）用于观察周围多个方向。红外相机（infrared camera）记录红外波段的信号。雷达（radar）利用电磁波回波，X 射线成像（X-ray imaging）利用穿透后的测量，计算机断层成像（computed tomography, CT）从多方向测量重建内部结构。测距相机（range camera）可记录距离，深度图（depth map）的像素值代表远近，不一定代表明暗。

因此，“图像里的亮点”要结合数据含义解释。普通灰度照片里的高值通常显示得较亮；深度图中的高值可能只是更远的地方。增强或识别之前，必须知道像素记录的是什么量。

原课程的应用案例可以沿着“观测对象、希望得到的答案”来理解。门禁与客流统计要知道谁进出、多少人经过；A--B 门防尾随系统要判断两道门之间是否多出一个人；遗留异物检测要发现原来没有而持续出现的物体；车牌识别要把字符图像转成文字；机场飞机起降跟踪要估计飞机位置随时间怎样变化；无人机航拍可用于交通参数统计和竣工图测绘。它们对分辨率、光照、时间采样和几何校正的要求不同，因此不能只靠一种“万能增强”。

\subsection*{原题 F2 低级与高级图像处理}
\textbf{题目。}低级图像处理与高级图像处理的区别？本课程的重点是哪一类？

\textbf{完整解答。}先比较输出。低级处理常把输入图像变成另一幅图像，目的是改善可见性或数据质量，如去噪、平滑、锐化和对比度调整；高级处理更关注从图像中获取结构或语义，如分割对象、描述特征、识别类别。再比较服务对象：原课程把前者重点描述为面向人的观察，把后者重点描述为面向机器的自主感知；实际系统中，增强也经常服务机器识别。最后指出课程范围：本课程重点是低级图像处理。作答时给出“去噪后仍是一张图”与“识别后得到车牌号”两个例子，就能说明区别，而不只是重复名词。

\section{人眼为什么会改变对同一张图的判断}
\label{sec:foundations-eye}

从室外走进较暗的房间，最初可能看不清家具，停留一段时间后细节逐渐出现。家具没有突然变亮，变化的是视觉系统的工作状态。图像增强常是为人观看服务的，因此仅知道像素数值还不够，还要理解人眼怎样接收光，以及周围环境怎样影响明暗判断。

\subsection{光怎样到达视网膜}

角膜（cornea）是眼球前方透明、参与折射的结构；巩膜（sclera）是外部的坚韧保护层。虹膜（iris）中央的开口是瞳孔（pupil），其大小改变进入眼内的光量。晶状体（lens）进一步帮助光线聚焦。眼球后部的视网膜（retina）包含把光信号转成神经信号的感光细胞（photoreceptor）。脉络膜（choroid）位于巩膜与视网膜之间，参与供血并吸收杂散光。视神经（optic nerve）离开眼球的位置没有通常的感光细胞，形成盲点（blind spot）。

\begin{figure}[htbp]
\centering
\begin{tikzpicture}[x=1cm,y=1cm,font=\small]
\draw[thick] (0,0) ellipse (3 and 1.75);
\draw[thick] (-2.75,.65) .. controls (-3.55,.2) and (-3.55,-.2) .. (-2.75,-.65);
\draw[very thick] (2.3,1.08) arc (34:-34:2.78 and 1.92);
\draw[thick] (-1.85,0) ellipse (.3 and .85);
\draw[very thick] (-2.3,.95)--(-2.3,.35);
\draw[very thick] (-2.3,-.95)--(-2.3,-.35);
\draw[-{Latex},thin] (-4.4,.6)--(-2.2,.3)--(-1.6,.25)--(2.7,-.55);
\draw[-{Latex},thin] (-4.4,-.6)--(-2.2,-.3)--(-1.6,-.25)--(2.7,.55);
\draw[thick] (2.6,-.8)--(3.5,-1.2);\draw[thick] (2.8,-1.05)--(3.45,-1.45);
\node[anchor=east] at (-3.7,1.2) {角膜};\draw (-3.55,1.1)--(-3.05,.45);
\node at (-2.5,1.55) {虹膜};\draw (-2.5,1.35)--(-2.3,.75);
\node at (-.95,1.5) {晶状体};\draw (-1,1.25)--(-1.7,.65);
\node at (.8,2.05) {巩膜};\draw (.8,1.85)--(.8,1.68);
\node[anchor=west] at (3.25,.75) {视网膜};\draw (3.2,.65)--(2.78,.3);
\node[anchor=west] at (3.55,-1.3) {视神经};
\node[align=center] at (.3,-2.35) {示意图只表示位置与光路关系，未按真实解剖比例绘制};
\end{tikzpicture}
\caption{眼睛将进入的光线聚焦到视网膜上，再由神经系统处理。}
\end{figure}

晶状体与普通固定玻璃镜片不同，它可以改变形状。睫状体（ciliary body）及相关结构调节晶状体的形状，这种为不同距离重新聚焦的过程叫调节（accommodation）。一般而言，看近处时晶状体较凸，看远处时较扁。近视（myopia）与远视（hyperopia）涉及眼的屈光能力、眼轴长度和成像位置，并不只是“晶状体这一次变厚或变薄”这一动作的另一种叫法。

\subsection{视锥与视杆各自适合怎样的工作}

视锥细胞（cone）在较明亮环境中发挥主要作用，支持颜色感知和精细空间辨别。视网膜中央的中央凹（fovea）具有很高的视锥密度，所以把字放在视线正中通常更容易看清。视锥主导的状态叫明视觉（photopic vision）。

视杆细胞（rod）在较暗环境中更敏感，但不能提供与视锥相同的颜色信息和精细空间分辨能力。视杆分布于较广的视网膜区域，中心凹的中心缺乏视杆；视杆主导的状态叫暗视觉（scotopic vision）。从明到暗的过渡环境中，两类细胞可共同起作用，称为中间视觉（mesopic vision）。“暗光下视锥不工作”适合作为极暗条件下的简化说法，不应理解为光线稍暗，视锥就全部关闭。

\begin{center}
\begin{tabularx}{\linewidth}{@{}p{2.1cm}XX@{}}
\toprule
比较项目 & 视锥细胞 & 视杆细胞\\
\midrule
主要适用环境 & 较明亮，主导明视觉 & 昏暗，主导暗视觉\\
颜色信息 & 不同类型视锥共同支持颜色辨别 & 不提供正常的颜色辨别\\
细节与分布 & 中央凹密集，精细辨别能力强 & 分布较广，暗光敏感但精细辨别较弱\\
对图像设计的启发 & 良好照明有助于看清文字和细小边缘 & 很暗环境下不宜仅用颜色编码关键差别\\
\bottomrule
\end{tabularx}
\end{center}

这一区别解释了为什么夜间只用红绿两种颜色标记信息可能不够可靠，也解释了为什么“能看见某个亮点”不等于“能分辨亮点内部的细节”。但不能由此推出“亮度感觉主要总是由视杆完成”。白天、正常室内照明和常见屏幕观看条件下，视锥对亮度与细节判断都十分重要。

\subsection{亮度适应与可察觉的变化}

物理强度描述光的大小；感知亮度（perceived brightness）描述人觉得有多亮。两者相关，但不是简单的一对一数字复制。视觉系统会根据环境调整灵敏度，这叫亮度适应（brightness adaptation）。当前适应的背景水平称为亮度适应级（brightness adaptation level）。

原课程用约 $10^{-6}$ 至 $10^4\ \mathrm{mL}$ 描述视觉系统跨不同适应状态可处理的很宽亮度范围。这里 $\mathrm{mL}$ 是毫朗伯（millilambert），不是体积单位毫升。这个约十个数量级的范围是经典教学示意，不能当作任何观察者、任何实验下都精确相同的端点。关键结论是：\textbf{眼睛不能在同一适应状态下同时清楚分辨这个巨大范围内的所有亮暗层次。}就像相机一次曝光不一定同时保留太阳附近和阴影深处的细节，人眼的瞬时辨别也受到当前状态限制。

现在考虑一个具体实验：背景亮度为 $I$，在背景上增加一块很小的亮度差 $\Delta I$。符号 $\Delta$ 读作“变化量”，所以亮块的亮度是 $I+\Delta I$。差别太小时，人可能无法稳定辨认；增加到某个水平后，就更容易察觉。用约定的实验判据测得的可察觉增量叫差别阈限（difference threshold）。原题把回答“看见了”的比例达到 $50\%$ 时的增量记为 $\Delta I_{50}$。下标 50 是实验判据的标记，不是把 $\Delta I$ 乘 50。

韦伯比（Weber ratio）定义为
\[
  W=\frac{\Delta I_{50}}{I},\qquad I>0.
\]
$I$ 与 $\Delta I_{50}$ 必须用相同的物理量、相同单位；分子分母单位相消，所以 $W$ 没有单位。$W=0.02$ 表示增量达到背景的 $2\%$ 时达到所用察觉判据。比值越小，说明在该实验条件下能够察觉的\emph{相对变化}越小，亮度分辨能力越好。

\begin{figure}[htbp]
\centering
\begin{tikzpicture}[font=\small,x=1cm,y=1cm]
\fill[gray!30] (0,0) rectangle (4.3,1.8);
\fill[gray!27] (1.6,.5) rectangle (2.7,1.3);
\draw (0,0) rectangle (4.3,1.8);
\node at (2.15,2.2) {背景 $I$ 与亮块 $I+\Delta I$};
\node[align=center] at (2.15,-.5) {比较的是差值相对于背景的比例\\图中灰度仅为示意，不模拟实验阈限};
\draw[-{Latex}] (5.1,0)--(11.1,0) node[right] {$W$};
\foreach \x/\lab in {5.1/0,7.1/{0.02},9.1/{0.04}} {\draw (\x,.08)--(\x,-.08);\node[below] at (\x,-.08) {$\lab$};}
\draw[thick] (7.1,.15)--(7.1,1.15);\node[above,align=center] at (7.1,1.15) {$2/100$\\相对阈限 $2\%$};
\draw[thick] (9.1,.15)--(9.1,1.15);\node[above,align=center] at (9.1,1.15) {$16/400$\\相对阈限 $4\%$};
\draw[-{Latex}] (10.2,-1)--(6,-1) node[midway,below] {越小，能察觉的相对变化越小};
\end{tikzpicture}
\caption{韦伯比必须用相对量比较。绝对增量变大，并不自动代表相对分辨能力变差。}
\end{figure}

在一定范围和固定实验条件下，差别阈限可能近似与背景强度成比例，即 $\Delta I_{50}\approx kI$。这叫韦伯定律（Weber's law）的近似形式；$k$ 是该条件下近似不变的比例。如果背景变为原来的 4 倍，阈限也恰好变成 4 倍，韦伯比不变。因此，“更亮背景需要更大的绝对增量”与“更亮背景下相对辨别更差”不能混为一谈。在非常暗的环境以及不同适应、观察时间、目标大小等条件下，这个简单近似也可能不成立。

\subsection*{原题 E4 韦伯比的计算与解释}
\textbf{题目。}（1）写出韦伯比的定义并说明其含义；（2）某背景亮度 $I=100$（单位任意但前后一致），在题设 $50\%$ 察觉判据下的增量 $\Delta I_{50}=2$，求韦伯比；（3）换成更亮背景 $I=400$、$\Delta I_{50}=16$，比较两次的亮度分辨能力。

\textbf{第一步，写定义。}$W=\Delta I_{50}/I$，表示达到该察觉判据所需的相对亮度变化。不是直接比较亮块总亮度 $I+\Delta I_{50}$。

\textbf{第二步，代入第一组。}
\[
 W_1=\frac{2}{100}=0.02=2\%.
\]
背景每 100 个强度单位增加 2 个单位，就达到题设判据。

\textbf{第三步，代入第二组。}
\[
 W_2=\frac{16}{400}=0.04=4\%.
\]
背景每 100 个强度单位需要增加相当于 4 个单位，才能达到同样判据。

\textbf{比较结论。}$W_1<W_2$，所以在这两组指定实验条件下，第一组的\emph{相对亮度分辨能力}较好。这个结论只比较题中两次测量，不能推出“环境越暗越容易分辨亮度”，也不能证明“暗适应总比明适应分辨能力强”。若第二组给出的是 $\Delta I_{50}=8$，则 $8/400=0.02$，两组相对阈限就相同。

\subsection{周围的像素会改变中心的视觉效果}

把完全相同的灰色小方块分别放到黑色和白色背景上，黑背景上的小方块往往显得更亮。这叫同时对比（simultaneous contrast）：某处的感知亮度不仅取决于它本身，还取决于周围区域。检查下图中心小块的数字可知，两块都用相同填充绘制；视觉上的不同不是计算机改了中间数值。

\begin{figure}[htbp]
\centering
\begin{tikzpicture}[font=\small,x=1cm,y=1cm]
\fill[black!85] (0,0) rectangle (4.3,2.2);
\fill[black!8] (5.1,0) rectangle (9.4,2.2);
\draw (0,0) rectangle (4.3,2.2);\draw (5.1,0) rectangle (9.4,2.2);
\fill[black!45] (1.55,.5) rectangle (2.75,1.7);
\fill[black!45] (6.65,.5) rectangle (7.85,1.7);
\node at (2.15,-.38) {暗背景};\node at (7.25,-.38) {亮背景};
\node at (4.7,-.9) {两个中心方块的实际填充灰度完全相同};
\end{tikzpicture}
\caption{同时对比的示意图。打印条件可能改变错觉强弱，但中心数值相同这一点不变。}
\end{figure}

马赫带（Mach bands）是另一种例子。在某些明暗过渡的交界附近，人会觉得存在额外的亮边或暗边，而物理亮度并没有那样的局部峰谷。视觉系统的侧抑制（lateral inhibition）等机制有助于解释这种边缘强化现象：一个位置的神经反应会受到邻近位置影响。视错觉（optical illusion）则还包括对长度、方向和形状的误判。它们共同提醒我们，图像数值的客观变化与观看者的主观判断不能简单画等号。

对图像增强来说，这意味着均匀分布的灰度数值不一定最好看，更亮也不一定更清楚。判断结果时要说明观察任务：看文字是否清晰，找裂缝是否容易，还是颜色是否自然。与此同时，主观“看上去舒服”不能代替测量精度或识别准确率。

\subsection*{原题 F3 视锥与视杆及其设计意义}
\textbf{题目。}视锥细胞与视杆细胞有什么差别？这对图像设计有什么意义？

\textbf{完整解答。}视锥在较明亮环境中支持颜色与精细细节辨别，中央凹密度高；视杆在较暗环境中更敏感、分布较广，但颜色辨别和精细空间分辨较弱。设计供人观看的图像时，应兼顾观察照明、文字与背景的对比，以及颜色之外的形状或明暗标记。评价增强效果，还要考虑亮度适应与同时对比等视觉现象。

不能把这个结论延伸成“中值滤波不会损失视觉信息”。中值滤波可能去掉小目标、细线和纹理，其效果必须结合目标与噪声检查；细胞的空间汇聚也不是中值滤波无损的证明。视锥、视杆的区别给出的是视觉设计依据，而不是任意图像算法都有效的保证。

\section{从连续场景到有限像素及像素之间的关系}
\label{sec:foundations-digitization}

拍照时，真实世界并不会预先站在一张整数网格上。空间位置可以连续变化，接收到的光强也可以连续变化。计算机却只能保存有限个位置和有限精度的数值。数字化（digitization）因此包含两个彼此独立的选择：在哪里测量，以及把测量结果记成哪些数。

\subsection{采样决定在哪里量 量化决定能记多细}

先用测温作类比。沿走廊每隔一米放一个温度计，是在选择测量位置；把 $23.37\,^{\circ}\mathrm{C}$ 记成 $23\,^{\circ}\mathrm{C}$ 或 $23.4\,^{\circ}\mathrm{C}$，是在选择数值精细程度。增加温度计个数，不会自动增加每支温度计的小数位；保留更多小数，也不会补回没有放温度计的地方。

采样（sampling）是在空间上离散化：从连续位置选出一组测量位置。对二维图像，可以每隔 $\Delta x$ 选一列、每隔 $\Delta y$ 选一行。$\Delta x$ 和 $\Delta y$ 是采样间隔（sampling interval），单位可以是毫米、微米，或经过成像几何换算后的物体表面距离。符号 $\Delta x$ 表示相邻采样位置之间的横向距离，不是灰度变化。

量化（quantization）是在幅度上离散化：把测得的连续强度归入有限个灰度级。假设一个设备把归一化强度从 0 到 1 分成四档，并用 0、1、2、3 编码，则强度 $0.10$ 和 $0.20$ 可能都进入第 0 档，强度 $0.68$ 进入第 2 档。归一化（normalization）在这里是把原范围按比例变成便于描述的 0 到 1，不是引入新的物理亮度单位。

\begin{figure}[htbp]
\centering
\begin{tikzpicture}[font=\small,x=1cm,y=1cm]
\begin{scope}
\draw[-{Latex}] (0,0)--(5.4,0) node[right]{位置};
\draw[-{Latex}] (0,0)--(0,2.4) node[above]{强度};
\draw[thick,domain=0:5,samples=120] plot (\x,{1.2+.75*sin(95*\x)});
\foreach \x in {.4,1.2,2,2.8,3.6,4.4}{
\pgfmathsetmacro{\yy}{1.2+.75*sin(95*\x)}
\draw[dashed,gray] (\x,0)--(\x,\yy);\fill (\x,\yy) circle (1.8pt);
}
\draw[<->] (.4,-.35)--(1.2,-.35) node[midway,below] {$\Delta x$};
\node at (2.5,2.85) {采样：只取选定位置的值};
\end{scope}
\begin{scope}[xshift=7.5cm]
\draw[-{Latex}] (0,0)--(4.5,0) node[right]{输入};
\draw[-{Latex}] (0,0)--(0,2.5) node[above]{输出级};
\foreach \y/\lab in {.25/0,.9/1,1.55/2,2.2/3}{\draw[dotted,gray] (0,\y)--(4,\y);\node[left] at (0,\y){\lab};}
\draw[very thick] (0,.25)--(1,.25)--(1,.9)--(2,.9)--(2,1.55)--(3,1.55)--(3,2.2)--(4,2.2);
\foreach \x/\lab in {0/0,1/{.25},2/{.50},3/{.75},4/1}{\draw (\x,0)--(\x,-.06);\node[below] at (\x,-.06){$\lab$};}
\node at (2,2.85) {量化：一段输入归为同一级};
\end{scope}
\end{tikzpicture}
\caption{采样改变位置的密度；量化改变允许输出的数值。两者不能互相替代。}
\end{figure}

上图右侧是一种四档量化规则，规定 $[0,0.25)$ 对应 0，$[0.25,0.50)$ 对应 1，依此类推，端点 1 归最后一档。方括号表示包含端点，圆括号表示不包含端点。不同设备也可以使用“舍入到最近代表值”的规则；做数值题时应说明采用哪一种，不能把分档编号与该档的物理代表值混淆。

若以固定步长 $q$ 对连续值做最近值舍入，在不超过量化器范围的条件下，量化误差（quantization error）的绝对值至多为 $q/2$。例如允许值为 $0,0.1,0.2,\ldots$，输入 $0.26$ 舍入为 $0.3$，误差为 $0.04$，没有超过 $0.05$。这里“误差”是存下来的代表值与原值之差，不是某种额外随机噪声必然服从的分布。

实际传感器的像素会对一个小面积、一个曝光时段内的光作响应，并不真是无限小点上的瞬时测量。把它当作网格点采样，是便于学习和计算的模型。像素有限面积、镜头模糊与曝光运动都会影响最终细节，后面讨论分辨率时不能忽略这些因素。

更完整的光场可抽象写成 $I=E(x,y,z,\lambda,t)$。这句话只表示光的测量值可能依赖空间位置 $(x,y,z)$、波长 $\lambda$ 和时间 $t$；$E$ 是描述这种依赖关系的函数名称，不要求在这里进行任何微积分运算。固定观察平面、某个波段和某段拍摄时间后，才得到常用的二维图像 $f(x,y)$。

\subsection{空间分辨率和灰度分辨率不能混为一谈}

空间分辨率（spatial resolution）关注空间上多小的结构还能被区分。数字网格的行列数是相关因素，但必须连同覆盖范围理解：同样 1000 列，覆盖 10 毫米与覆盖 10 千米的实际细节能力完全不同。在同样视场内，更多有效像素通常意味着更细的空间采样；是否真正分辨出来，还受镜头、对焦、运动模糊、噪声和目标对比度影响。

灰度分辨率（gray-level resolution）关注可区分的强度档数，常用每像素位数 $b$ 或级数 $L=2^b$ 描述。空间采样变稀，斜边容易呈锯齿，小目标容易消失；灰度级太少，一段平滑的明暗变化可能变成一条条可见的带状层次，称为假轮廓（false contour）或色带（banding）。

\begin{figure}[htbp]
\centering
\begin{tikzpicture}[font=\small,x=1cm,y=1cm]
\foreach \x in {0,...,3}{\foreach \y in {0,...,3}{
\pgfmathtruncatemacro{\shade}{15+20*\x}
\fill[black!\shade] (.55*\x,.55*\y) rectangle ++(.55,.55);
}}
\draw[step=.55,white] (0,0) grid (2.2,2.2);\draw (0,0) rectangle (2.2,2.2);
\node[align=center] at (1.1,-.6) {位置稀疏\\小细节难以记录};
\foreach \x in {0,...,15}{\foreach \y in {0,...,15}{
\pgfmathtruncatemacro{\shade}{15+5*\x}
\fill[black!\shade] (4+.1375*\x,.1375*\y) rectangle ++(.1375,.1375);
}}
\draw (4,0) rectangle (6.2,2.2);
\node[align=center] at (5.1,-.6) {位置更密\\灰度也较丰富};
\foreach \x in {0,...,3}{\pgfmathtruncatemacro{\shade}{15+20*\x}\fill[black!\shade] (8+.55*\x,0) rectangle ++(.55,2.2);}
\draw (8,0) rectangle (10.2,2.2);
\node[align=center] at (9.1,-.6) {灰度只有四档\\渐变出现分带};
\end{tikzpicture}
\caption{采样不足与灰度档数不足产生不同的损失。网格线与条带为示意，不能据此量化真实设备的分辨率。}
\end{figure}

\subsection*{原题 F4 采样和量化}
\textbf{题目。}什么是采样？什么是量化？图像数字化包含哪两方面？

\textbf{完整解答。}采样把连续空间坐标变成有限网格位置，决定在哪里取值；量化把连续的强度变成有限个灰度级，决定每个位置的测量值怎样编码。数字化包含这两个方面。例如，每隔 $0.1\,\mathrm{mm}$ 测一次是采样选择，把每次测量存成 0 到 255 的整数是量化选择。前者的间隔单位是空间长度，后者的档数是数量、位数单位是 bit。写出这一具体例子，就能避免把“增加像素”与“增加灰度级”当作同一件事。

\subsection*{原题 F5 两种分辨率}
\textbf{题目。}空间分辨率与灰度分辨率的区别？分别由什么决定？

\textbf{完整解答。}空间分辨率描述区分空间细节的能力；在固定视场和其他成像条件下，像素采样越密，能够记录的空间变化通常越细。灰度分辨率描述数值上可用的强度级数，$b$ 比特对应 $2^b$ 个等级。前者不足常表现为锯齿、块状和细节丢失，后者不足常表现为渐变分带、假轮廓。必须补充条件：图像宽高只能描述网格规模，不能脱离物方覆盖范围、镜头与成像质量单独保证实际分辨率；把模糊图扩大成更多像素也不等于获得更多真实细节。

\subsection{周期和频率描述条纹有多密}

考虑一排重复的亮暗条纹。从一条亮纹的中心走到下一条相同亮纹的中心，是一个完整空间周期（spatial period），记为 $T$。若每隔 $0.2\,\mathrm{mm}$ 重复一次，则 $T=0.2\,\mathrm{mm}$。空间频率（spatial frequency）是单位长度内重复的次数，记为 $f$：
\[
  f=\frac{1}{T}.
\]
所以 $T=0.2\,\mathrm{mm}$ 时，$f=5\ \mathrm{cycles/mm}$，意思是每毫米有 5 个完整周期。周期越短，频率越高。“高频”不是说像素值很大，而是说灰度随位置变化得快。

采样频率（sampling frequency）同样是单位长度内取样的次数，记为 $f_s$。若相邻采样点间隔为 $\Delta x$，则
\[
  f_s=\frac{1}{\Delta x}.
\]
例如 $\Delta x=0.05\,\mathrm{mm}$，每毫米取 $1/0.05=20$ 次样，故 $f_s=20\ \mathrm{samples/mm}$。\emph{cycles/mm} 数的是信号的重复周期，\emph{samples/mm} 数的是测量点；两者有相同的“每长度”维度，但不是同一个对象。

\subsection{采样太稀为什么会产生假的条纹}

每隔很长距离才测一次条纹，就可能错过中间的明暗变化。更严重的是，稀疏测得的一串值可能同时符合两种不同条纹：原来的快速变化，以及另一种较慢变化。数字系统无法仅凭这串值判断哪种是真的，这叫混叠（aliasing）。相机拍摄细密织物、显示器或栅栏时出现的粗条纹，常与这种现象有关。

奈奎斯特采样定理（Nyquist sampling theorem）讨论一个理想条件：信号的频率成分存在最高上限 $f_{\max}$，超过它的成分完全不存在，这叫带限（band-limited）。在理想采样与理想重建等条件下，为避免最高频率正好落在含糊的边界，采用
\[
  f_s>2f_{\max},\qquad
  \Delta x<\frac{1}{2f_{\max}}.
\]
其中 $f_{\max}$ 是信号中最高的空间频率，$f_s$ 是采样频率，$\Delta x$ 是采样间隔。把间隔放在分母会使不等号方向反过来：要求频率更高，等价于要求间隔更小。

一些教材把临界条件写成 $f_s\geq2f_{\max}$，方便记忆“每周期至少两个样本”。但是等号是理想边界，不应作为对任意相位都可靠的工程设计点。相位（phase）在这里先理解为条纹相对采样网格平移了多少：同一排条纹平移半格，亮暗峰位置就会落到不同采样位置。

\begin{figure}[htbp]
\centering
\begin{tikzpicture}[font=\small,x=1.65cm,y=1cm]
\draw[-{Latex}] (0,0)--(6.2,0) node[right]{位置 $x$};
\draw[-{Latex}] (0,-1.3)--(0,1.5) node[above]{强度变化};
\draw[thick,domain=0:6,samples=180] plot (\x,{sin(180*\x)});
\foreach \x in {0,...,6}{\fill (\x,0) circle (2pt);\draw[dotted] (\x,-1.2)--(\x,1.1);}
\draw[very thick,dashed] (0,0)--(6,0);
\draw[<->] (0,-1.6)--(2,-1.6) node[midway,below]{一个周期 $T$};
\node[align=center] at (4,1.65) {每隔 $T/2$ 采样一次\\恰好全落在零点};
\node[align=center] at (4,-1.65) {采样结果与一条恒定线相同\\真实的振荡已无法由样本判定};
\end{tikzpicture}
\caption{临界采样的相位风险。曲线表示相对某背景的亮度变化，允许出现负的变化量，并非负的实际光强。}
\end{figure}

上图的真实曲线在采样点之间仍然起伏，但所有记录值都相同。增加一个足够大的背景亮度，就得到完全非负的真实强度，问题仍然存在。因此这不是“实际图像不会有负值”可以排除的数学小把戏，而是临界位置确实可能漏掉变化。

真实图像中的尖锐边缘并不严格带限。实际系统通常先用光学或数字低通来削弱过细变化，再采样；这叫抗混叠（anti-aliasing）。降低图像尺寸时也一样，不能只随手每隔几行几列丢掉像素，否则新的采样网格可能把细纹理误认为粗纹理。先适当平滑会牺牲一部分细节，却能减轻虚假结构。

\subsection*{原题 E5 从最细周期计算采样要求}
\textbf{题目。}某图像最细结构的空间周期为 $0.2\,\mathrm{mm}$，要在数字化时不发生混叠，采样间隔最大为多少？采样频率至少为多少？

\textbf{第一步，明确模型。}按题目的理想周期模型，把最细周期 $T_{\min}=0.2\,\mathrm{mm}$ 对应的频率作为最高频率。若所谓“结构”是尖锐条纹，还包含更高频成分，则需要先说明带限或预滤波假设，不能把条纹宽度机械代入这个周期公式。

\textbf{第二步，周期变频率。}
\[
 f_{\max}=\frac{1}{T_{\min}}
 =\frac{1}{0.2\,\mathrm{mm}}
 =5\ \mathrm{cycles/mm}.
\]
\textbf{第三步，求临界值。}
\[
 2f_{\max}=10\ \mathrm{samples/mm},\qquad
 \frac{1}{2f_{\max}}=0.1\,\mathrm{mm}.
\]
因此，原讲义习惯写的答案是“至少 10 次/mm、间隔不大于 $0.1\,\mathrm{mm}$”。更严谨地说，$10$ 次/mm 与 $0.1\,\mathrm{mm}$ 是\textbf{临界界限}：若要求对一般相位可靠，应取 $f_s>10$ 次/mm、$\Delta x<0.1\,\mathrm{mm}$，而不把等号当作安全点。

\textbf{第四步，给出有裕量的例子。}取 $\Delta x=0.05\,\mathrm{mm}$，则 $f_s=20$ 次/mm，一个最细周期内有 $0.2/0.05=4$ 个采样间隔。取 $0.04\,\mathrm{mm}$ 则对应 25 次/mm。它们比临界条件更稳妥，但“选四点或五点”仍不是任何镜头、噪声和目标条件下的质量保证，真实系统还需成像测试。

\subsection{视场把传感器尺寸换成物体上的尺寸}

相机参数中的像元尺寸（pixel pitch）描述传感器上相邻采样单元的间距，例如 $4.52\,\mu\mathrm{m}$。拍到的物体往往比传感器大得多，因此传感器上一个像元不等于物体上同样大小的一块区域。物方像素尺寸（object-space pixel size），在遥感中常称地面采样距离（ground sample distance, GSD），才描述相邻像素中心在被拍平面上相隔多远。

视场（field of view, FOV）在本节指相机在某个物体平面上覆盖的实际宽或高，不是镜头的视角角度。焦距（focal length）$f$ 描述成像几何的一项镜头参数；物距（object distance）$Z$ 描述物体平面离相机的距离。对近似正对相机的平面、较远物体，并忽略畸变，可以用相似三角形得到
\[
 \frac{W}{S_x}\approx\frac{Z}{f},
 \qquad W\approx S_x\frac{Z}{f}.
\]
式中 $W$ 是横向物方视场，$S_x$ 是传感器有效宽度；同理 $H\approx S_yZ/f$。相似三角形（similar triangles）只是说两只三角形形状相同，对应边按相同比例放大。物距远大于焦距时，像距可以近似用焦距代替；近距离精密测量时应使用更准确的成像模型并做标定。

\begin{figure}[htbp]
\centering
\begin{tikzpicture}[font=\small,x=1cm,y=1cm]
\draw[thick] (0,-1.8)--(0,1.8);
\draw[thick] (8,-.65)--(8,.65);
\fill (6,0) circle (2pt);
\draw (0,1.8)--(6,0)--(8,-.6);
\draw (0,-1.8)--(6,0)--(8,.6);
\draw[dashed] (-.5,0)--(8.6,0);
\draw[<->] (-.5,-1.8)--(-.5,1.8) node[midway,left] {$W$};
\draw[<->] (8.5,-.6)--(8.5,.6) node[midway,right] {$S_x$};
\draw[<->] (0,-2.25)--(6,-2.25) node[midway,below] {物距 $Z$};
\draw[<->] (6,-1.2)--(8,-1.2) node[midway,below] {约为 $f$};
\node at (0,2.15) {物体平面};\node at (6,.5) {成像中心};\node at (8,1.05) {传感器};
\node[align=center] at (4,-3.1) {图为放大的几何示意；实际物距通常远大于焦距};
\end{tikzpicture}
\caption{相似三角形把传感器上的长度换算为物体平面上的长度。}
\label{fig:foundation-fov}
\end{figure}

若横向有 $N_x$ 个像素、像元间距为 $p$，则常用有效宽度近似为 $S_x=N_xp$。这个表达把每个像素看作宽度 $p$ 的覆盖单元；若只讨论第一个到最后一个像素\emph{中心}的距离，则是 $(N_x-1)p$，两种口径应按需求区分。在本题的视场估算中使用前一种常见工程口径。代入可得
\[
 W\approx N_xp\frac{Z}{f},\qquad
 d\approx\frac{W}{N_x}=p\frac{Z}{f}.
\]
$d$ 是一个像素在物方对应的长度。例如 $d=0.2\,\mathrm{mm/pixel}$，表示每像素覆盖的标称物方宽度约为 $0.2\,\mathrm{mm}$；10 个像素约对应 2 毫米。这里的 pixel 是计数说明，不是另一个物理长度单位。

公式也展示出直观关系。在其他条件不变时，拍摄更远的物体会覆盖更大的视场，物方像素也更大；增加焦距会缩小视场，物方采样变细。只增加传感器像素数而保持同样像元尺寸、焦距和物距，主要扩大覆盖范围，并不自动减小 $d$。固定视场下增加像素数，才直接把 $d=W/N_x$ 变小。

\subsection*{原题 E2 分辨率的工程换算}
\textbf{题目。}某相机像素数 $7952\times5304$，像素尺寸 $4.52\,\mu\mathrm{m}$，物距 $300\,\mathrm{m}$，镜头焦距 $36.18\,\mathrm{mm}$。求：（1）传感器尺寸；（2）视场，横向乘纵向；（3）图像中一个像素代表多大的实际尺寸。

\textbf{第一步，整理单位。}
\[
 1\,\mathrm{mm}=1000\,\mu\mathrm{m},\quad
 1\,\mathrm{m}=1000\,\mathrm{mm},\quad
 p=4.52\,\mu\mathrm{m}=0.00452\,\mathrm{mm}.
\]
将像元尺寸先换成毫米，可以直接与焦距的毫米相消。像素数是计数，不需要再换长度单位。

\textbf{第二步，求传感器有效尺寸。}
\begin{align*}
 S_x&=7952\times0.00452\,\mathrm{mm}=35.94304\,\mathrm{mm},\\
 S_y&=5304\times0.00452\,\mathrm{mm}=23.97408\,\mathrm{mm}.
\end{align*}
因此传感器约为 $35.94\,\mathrm{mm}\times23.97\,\mathrm{mm}$。保留中间计算的足够位数，可以减少反复四舍五入造成的差异。

\textbf{第三步，求物方视场。}在上述远距离、正视平面的近似下，
\begin{align*}
 W&=35.94304\,\mathrm{mm}\times
       \frac{300\,\mathrm{m}}{36.18\,\mathrm{mm}}
   \approx298.04\,\mathrm{m},\\
 H&=23.97408\,\mathrm{mm}\times
       \frac{300\,\mathrm{m}}{36.18\,\mathrm{mm}}
   \approx198.79\,\mathrm{m}.
\end{align*}
分子与分母中的毫米相消，最后留下米；若算出来仍是“毫米乘米”，就说明单位约分没有完成。

\textbf{第四步，求单像素的物方长度。}
\begin{align*}
 d&=0.00452\,\mathrm{mm}\times
       \frac{300\,\mathrm{m}}{36.18\,\mathrm{mm}}\\
  &\approx0.0374793\,\mathrm{m/pixel}
   =3.74793\,\mathrm{cm/pixel}
   \approx3.75\,\mathrm{cm/pixel}.
\end{align*}
也可以用 $298.04/7952$ 检查横向结果，用 $198.79/5304$ 检查纵向结果。方形像元和同一成像几何应给出相同的物方尺度。

\textbf{结果解释。}该参数组合的视场约为 $298.04\,\mathrm{m}\times198.79\,\mathrm{m}$，标称采样尺度约为 $3.75\,\mathrm{cm/pixel}$。原稿中的 $4.52\,\mathrm{cm/pixel}$ 与题设不相容，不能因为量级接近就当作同一个结果。物方长度明显小于一个像素的毫米级裂缝，通常不能靠这组参数可靠分辨其宽度；增加软件显示倍率不会把缺失测量补回来。
